\section{结论}
\label{sec:conclusion}

本文呈现一项关于可见度退化下知识蒸馏失效的机制驱动实证研究。通过15个控制实验的结构化分支级分析，我们剖析了现有KD方法在雾天条件下失效的原因，并识别了哪些假设机制得到实证支持。

我们的核心贡献是将问题从\emph{如何在雾天改进KD}重新框架为\emph{为什么KD在雾天下失效}。这一转变使得系统性机制测试成为可能，而非增量式方法开发。

我们的主要发现包括：

\textbf{第一}，分支级性能结构确实存在且显著。不同KD机制对可见度退化的响应不同，定位蒸馏表现出相对鲁棒性，而logit蒸馏尽管文献中广泛使用却被证明并非最优。

\textbf{第二}，两种常被假设的机制——分布错配和不确定性放大——与KD增益的相关性弱（$|r|<0.5$，$p>0.7$），在本设置中不足以解释观测到的退化。这些结果提示，关于KD失效的普遍直觉可能不完整。

\textbf{第三}，在所测试机制中，遮挡——可见度退化的具体可测成分——与定位性能呈现最强负相关（$r=-0.989$，$p=0.0015$），是最直接得到支持的机制。当遮挡隐藏物体边界时，准确定位所需的空间信息变得不可用。我们具体识别遮挡为可见度退化中最强得到支持的因素，同时指出可见度损失的其他方面（对比度下降、颜色偏移）可能贡献但未在本研究中直接测试。

\textbf{第四}，我们的分析表明logit KD失效不太可能仅通过超参数调整解决，提示需要架构层面而非配置层面的解决方案。

这些发现具有方法论意义：针对恶劣条件的未来KD研究应优先处理信息恢复和遮挡感知设计，而非仅关注分布对齐。主导干净域KD研究的机制可能无法迁移到可见度退化设置。

本研究也例证了更广泛的方法：机制驱动实证分析即使在产生负面结果时也能推进理解。排除不足的解释——分布错配、不确定性放大和超参数不足——本身就是进展，缩小了可行解决方案的搜索空间。

未来工作应将这些分析扩展到真实世界退化，验证遮挡感知KD架构，并探索多教师或自适应方法，以在严重可见度损失下维持有效性。目标不仅是提高基准分数，而是开发对表征恶劣视觉条件的信息损失具有根本鲁棒性的KD方法。

\section*{数据与代码可用性}

实验数据和分析脚本可在\url{https://github.com/godzhiwzz-create/kd-visibility-study}获取。Foggy Cityscapes数据集公开可获取于\url{https://www.cityscapes-dataset.com/}。

\section*{利益冲突声明}

作者声明没有已知的可能影响本文报告工作的竞争性财务利益或个人关系。
